\documentclass[11pt]{article}
\usepackage[a4paper,margin=25mm]{geometry}
\usepackage{amsmath,amssymb,amsthm,booktabs,hyperref}
\newtheorem{theorem}{Theorem}
\newtheorem{lemma}[theorem]{Lemma}
\newtheorem{corollary}[theorem]{Corollary}
\newcommand{\F}{\mathbb F}
\newcommand{\gen}[1]{\langle #1\rangle}
\newcommand{\CG}{\operatorname{C}}
\title{A polynomial prefix obstruction in the natural $G_2(11)$ model}
\author{human\_compression}
\date{31 October 2026}
\begin{document}
\maketitle

\begin{abstract}
An omitted involution which is a polynomial in an adjacent product belongs
to every geometric maximal overgroup of that product. In the natural
$G_2(11)$ representation, the remaining maximal types have no element of
order $33$ or $55$. This gives a conditional exclusion of an entire prefix
of rank-four markings. The finite missing premise is a single matrix
polynomial identity for each original prefix, independent of the fourth
mark. The accompanying finite certificates supply ten such identities; their verification is separate from the imported group classification. The argument uses
the retained maximal-subgroup classification and its natural geometric
identifications; it does not prove a new classification or an open
conjecture.
\end{abstract}

\section{The representation and the original obligation}

Let $C$ be a split Cayley algebra over $\F=\F_{11}$ and let $V=C_0$ be its
seven-dimensional pure space. We use the admitted full group
$G=\operatorname{Aut}_{\F}(C)\cong G_2(11)$ in its faithful natural action
on $V$. All changes of coordinates and transfers to the original marked
group must retain the previously acquired full-group, faithful-model,
original-label, member, orbit and transporter premises. Preservation of a
quadratic form alone does not establish this representation premise.

Fix nonidentity involutions $a,b,e\in G$, with $[a,e]=[b,e]=1$, and put
$r=ab$. For a proposed fourth mark $c\in G$, the intended ordered marking
is $(a,b,c,e)$ and its left end is $L=\gen{a,b,c}$. A complete original
classification must retain positive $G$-generating string C-groups and
exclude every other original marking. In particular, the original
$q=11$ result is a positive classification as well as an exhaustion; it
does not assert that all rank-four markings are absent.

The theorem below excludes only the specified prefixes. Its two branches
are deliberately distinct: a proper generated quartet fails generation of
$G$, while a left end equal to $G$ fails the intersection property. The
proof does not decide which branch occurs for each $c$.

\section{Polynomial membership and geometric overgroups}

\begin{lemma}\label{poly}
Let $W$ be a finite-dimensional vector space over a field $k$, let
$r\in\operatorname{End}_k(W)$, and let $e=f(r)$ for $f\in k[X]$.
Then every $r$-invariant subspace is $e$-invariant. Moreover $e$ commutes
with every endomorphism commuting with $r$.
\end{lemma}
\begin{proof}
If $U$ is $r$-invariant, every $r^i$ sends $U$ into $U$, and hence so does
their linear combination $f(r)$. If $t r=r t$, induction gives
$t r^i=r^i t$ for every $i\geq0$; linear combination gives $te=et$.
When $e$ is invertible, the inclusion $e(U)\subseteq U$ is an equality,
since its restriction to the finite-dimensional subspace is injective.
\end{proof}

We use the following exact geometric facts in the same natural action.
They are retained classical inputs, not consequences of abstract subgroup
names.

\begin{enumerate}
\item Each maximal parabolic is the full stabilizer of an isotropic pure
line, or of a null two-space on which Cayley multiplication vanishes.
\item The split and nonsplit quadratic-line maximal subgroups are the
full setwise stabilizers of their respective nonisotropic pure lines.
Their structures are $\operatorname{SL}_3(11):2$ and
$\operatorname{SU}_3(11):2$. These are line stabilizers: the vector
stabilizers are their index-two kernels and are not substituted for them.
\item Each maximal subgroup of involution-centralizer type is
$\CG_G(t)$ for an actual nonidentity involution $t\in G$.
\end{enumerate}

Wilson, Sections 4.2.5--4.2.6 of the retained preliminary edition, gives
the two parabolics as the indicated full one- and two-space stabilizers,
the two full line stabilizers, and the involution centralizer. In that
edition these are printed pages 105--108, physical pages 119--122.
The subspaces in (1) and (2) belong to the pure seven-space. Thus the
polynomial identity is required only on $V$. We do not assert that the
same polynomial also takes the value $1$ on the scalar summand of $C$.

\begin{lemma}\label{containers}
Suppose $a,b,e\in G$ and $e\in\F[r]$, where $r=ab$. Every geometric
maximal subgroup in the three preceding classes which contains $a,b$
also contains $e$.
\end{lemma}
\begin{proof}
Let $M$ contain $a,b$, so $r\in M$. If $M$ is the full stabilizer of
one of the indicated subspaces $U\subset V$, then $r(U)=U$.
Lemma~\ref{poly} gives $e(U)=U$. Since $e\in G$ and $M$ is the
\emph{full} stabilizer within $G$, this proves $e\in M$. The argument
applies to the setwise stabilizer of a line; it does not require its
vectors to be fixed rather than negated.

If $M=\CG_G(t)$, then $tr=rt$. Lemma~\ref{poly} gives $te=et$, and
therefore $e\in\CG_G(t)=M$. This step does not infer centralizer
membership from an arbitrary invariant eigenspace.
\end{proof}

\section{The exact nongeometric alternatives}

The retained complete maximal cover at $q=11$ has the two parabolics,
the involution-centralizer type, the two quadratic-line types, and the
four remaining types
\[
 2^3.\operatorname{PSL}_3(2),\qquad
 G_2(2),\qquad \operatorname{PGL}_2(11),\qquad J_1.
\]
The full list and the exceptional prime-field restrictions are the
previously reviewed input \emph{Short positive generation certificates
in the acquired $G_2(11)$ model}, SHA-256 prefix
\texttt{f1b6c8a69780c7b8}. In particular there is no proper subfield
case at the prime $11$, and the possible $\operatorname{PSL}_2(13)$
and $\operatorname{PSL}_2(8)$ exceptional cases are excluded by the
recorded field restrictions. They are not silently discarded by an
irreducibility assumption.

\begin{lemma}\label{exceptional}
None of these four remaining types contains an element of exact order
$33$ or $55$.
\end{lemma}
\begin{proof}
The orders of $2^3.\operatorname{PSL}_3(2)$ and $G_2(2)$ involve only
the primes $2,3,7$, and hence not $11$. In $\operatorname{PGL}_2(11)$
an element is semisimple of order dividing $10$ or $12$, or unipotent
of order $11$; neither $33$ nor $55$ occurs. The retained attributed
Janko spectrum has maximal orders under divisibility
\[
 \mu(J_1)=\{6,7,10,11,15,19\}.
\]
Neither $33$ nor $55$ divides a member of this set. This is the same
primary spectrum input already used by the accepted prefix-spectrum
argument: Li--Moghaddamfar--Vasil'ev--Wang,
arXiv:2312.07907v1, Lemma 3.9(2), retained through the
\emph{moved19/order4} note. Divisibility of $|J_1|$ alone would not
exclude the two orders and is not used in place of that spectrum.
\end{proof}

\section{The prefix obstruction}

\begin{theorem}\label{main}
Under the full natural-model premises above, suppose
\[
 |ab|\in\{33,55\},\qquad e\in\F[ab].
\]
For every $c\in G$, the ordered four-mark group $(a,b,c,e)$ is not a
string C-group generating $G$. This holds in particular for every
original selected $c$ satisfying the intended involution and string
conditions. No per-$c$ word, order, line or subgroup calculation is needed
for this exclusion.
\end{theorem}
\begin{proof}
Put $H=\gen{a,b,c,e}$ and $L=\gen{a,b,c}$. If $L<G$, choose a maximal
proper subgroup $M$ containing $L$. Its element $ab$ has exact order
$33$ or $55$. Lemma~\ref{exceptional} excludes all four nongeometric
types. Hence $M$ is one of the geometric types in
Lemma~\ref{containers}, which implies $e\in M$. Consequently
$H\leq M<G$ and the quartet does not generate the required group.

If $L=G$, then $e\in L$. Since $e\ne1$, the subgroups generated by
the disjoint label sets $\{0,1,2\}$ and $\{3\}$ have intersection
$\gen e\ne1$. This violates the intersection property. The two cases
exhaust the possibilities for $L$.
\end{proof}

This is a disjunctive exclusion of a target $G$-string C-group. It must
not be recorded as an unconditional IP failure for every $c$, nor as
unconditional nongeneration for every $c$. A typed adoption may retain
this logical alternative directly; if an existing ledger insists on
separate categories, its schema must first acquire an honest additional
category or the separate branches must be certified. The theorem does
not authorize relabelling old evidence.

\begin{corollary}[A sufficient cyclic certificate]\label{cyclic}
Suppose $[e,r]=1$ and $r$ is cyclic on $V$, meaning that some $v$ has
$v,vr,\ldots,vr^6$ as a basis. Then $e\in\F[r]$.
\end{corollary}
\begin{proof}
There is a polynomial $f$ of degree at most six with $ve=vf(r)$.
Commutation gives $(vr^i)e=(ve)r^i=vf(r)r^i=(vr^i)f(r)$.
The equality on the stated basis implies $e=f(r)$ on $V$.
\end{proof}

We do not infer cyclicity from an order alone, or from a scalar group
commutant. The accompanying finite proposal tests polynomial membership
directly. Cayley--Hamilton ensures that degree at most six suffices:
every higher power of a seven-dimensional operator is a linear
combination of $I,r,\ldots,r^6$.

\section{Finite certificate and scope}

For each exact original prefix, a positive carrier consists of seven
field elements $(u_0,\ldots,u_6)$ and the identity
\[
 e=\sum_{i=0}^{6}u_i(ab)^i.
\]
There are precisely 49 scalar entry equalities. The checker may propose
the coefficients by bounded elimination in the $49$ by $7$ linear
system, but must verify the full identity independently of that proposal
routine before returning a positive. An inconsistent system, missing
coefficient carrier, guard stop or failed check grants no prefix credit.
The actual order is checked by $r^m=I$ and
$r^{m/p}\ne I$ for each prime $p\mid m$, with $m=33$ or $55$.

The finite certificate package records the ten exact original rank-four prefixes
\[
156,160,166,167,168,170,174,175,183,187.
\]
Their seven coefficients and three ordered matrices are retained literally.
For prefixes $168$ and $170$, the exact adjacent product order is $33$;
for the other eight it is $55$. The accompanying verifier checks all
49 entries by Horner evaluation and checks the annihilating and proper
prime-divisor powers by independent repeated multiplication. It performs
no coefficient search or row reduction.

The theorem covers every fourth mark in each admitted original prefix.
The recorded finite domains comprise 74,330 original positions: 1,030
previously accepted nongenerating positions and 73,300 positions covered
by the disjunctive argument. These numbers describe those ten local
partitions. They do not give a net-new global paid-position count, and
the two alternatives of the theorem remain unseparated.

The comparison target is Table~2, p.~11, of Malcolm Hoong Wai Chen,
\emph{Abstract regular polytopes of $G_2(q)$ in even characteristic},
\href{https://arxiv.org/abs/2609.25004v1}{arXiv:2609.25004v1},
21 September 2026~\cite{Chen}. That table records 105 rank-four classes
at $q=11$ up to isomorphism and duality. The ten corresponding prefix jobs
have been closed in the current register. This partial resolution supplies
neither the positive-class count nor a completed field classification.
It supplies no whole-manuscript recovery, conjecture solution, Lean
verification or measured compression ratio.

\section{Attribution and relation to retained work}

The full-A2 end-intersection argument proposed during this review is
already present in the retained split/unitary end notes. The separate
\emph{Common point stabilizers in the natural Cayley group} note also
already proves the complete point-pair obstruction, including the
isotropic Heisenberg cases. Those proofs are reused, not rediscovered.
For proper ends, containment in their A2 containers does not imply that
their intersection contains the intersection of the containers.

The polynomial-invariant-subspace lemma, Cayley--Hamilton, the cyclic
commutant argument and the C-group definition are elementary standard
facts. The contribution here is their explicit common application to the
already acquired prefix model and the complete attributed $q=11$ cover.
No historical novelty is asserted. All classification, representation,
subgroup-spectrum, original-domain and model costs remain imported and
must be counted in any later compression audit.
\begin{thebibliography}{9}
\bibitem{Chen} Malcolm Hoong Wai Chen,
\emph{Abstract regular polytopes of $G_2(q)$ in even characteristic},
\href{https://arxiv.org/abs/2609.25004v1}{arXiv:2609.25004v1},
21 September 2026, Table~2, p.~11.
\bibitem{K} P.~B.~Kleidman,
The maximal subgroups of the Chevalley groups $G_2(q)$ with $q$ odd,
the Ree groups ${}^2G_2(q)$, and their automorphism groups,
\emph{J. Algebra} \textbf{117} (1988), 30--71.
\href{https://doi.org/10.1016/0021-8693(88)90239-6}{doi:10.1016/0021-8693(88)90239-6}.
\bibitem{N} H.~N.~Nguyen,
\emph{Irreducible restrictions of Brauer characters of the Chevalley group
$G_2(q)$ to its proper subgroups},
\href{https://arxiv.org/pdf/0910.4994}{arXiv:0910.4994v1},
proof of Theorem~4.4, printed pp.~10--11.
The complete list preceding the degree-based exclusions is the input used here.
\bibitem{W} R.~A.~Wilson, \emph{The Finite Simple Groups}, retained
preliminary edition, Sections~4.2.5--4.2.6, Table~4.1, printed pp.~105--109.
The published edition renumbers the Cayley chapter as Section~4.3.
\url{https://webspace.maths.qmul.ac.uk/r.a.wilson/TFSG.html}.
\bibitem{L} Li--Moghaddamfar--Vasil'ev--Wang,
\href{https://arxiv.org/pdf/2312.07907v1}{arXiv:2312.07907v1},
Lemma~3.9(2), printed p.~6 (the stated $J_1$ spectrum).
\end{thebibliography}
\end{document}
